\section{Duality frames and the compensator}\label{sec:frames}

\subsection{Three sectors}

\begin{definition}[Sectors]\label{def:sectors}
A kernel with \(s(z)>0\) is of \emph{matter type}, with \(s(z)<0\) of \emph{tachyon type}, and with \(s(z)=0\) \emph{light-like}.
\end{definition}

In the electromagnetic reading \eqref{eq:invariants}, matter-type kernels with real \(s\) are electric-dominated and tachyon-type kernels
magnetic-dominated.  The assignment of the names is a choice; the statements below depend only on the sign of \(s\).  The duality rotation
\(z\mapsto e^{i\alpha}z\) multiplies \(s\) by \(e^{2i\alpha}\).  In particular the quarter turn
\begin{equation}\label{eq:quarterturn}
z\mapsto iz,\qquad s\mapsto-s,
\end{equation}
exchanges the two massive sectors and fixes the null cone.  We use it as the frame change from matter to tachyon.

\begin{proposition}[Duality angle]\label{prop:duality}
Every \(z\) with \(s(z)\neq0\) can be written as \(z=e^{i\alpha}z_m\) with \(s(z_m)=|s(z)|>0\); \(\alpha=\tfrac12\arg s(z)\) is unique modulo \(\pi\) and
\(z_m\) unique up to sign.  Under \(SO(3,\mathbb C)\times U(1)_{\rm dual}\), acting by \(z\mapsto e^{i\alpha}Rz\), the orbits in
\(\mathbb C^3\setminus\{0\}\) are the shells \(\{|s|=r\}\), \(r>0\), and the punctured null cone.
\end{proposition}

\begin{proof}
\(s(e^{i\alpha}z_m)=e^{2i\alpha}s(z_m)\) is real positive exactly when \(2\alpha\equiv\arg s\pmod{2\pi}\).  The orbit statement follows from
Proposition~\ref{prop:orbits}, since the duality circle runs through all values of \(s\) of fixed modulus and fixes \(s=0\).
\end{proof}

So matter-type and tachyon-type kernels of equal \(|s|\) lie in one orbit, related by a quarter turn, a non-real \(s\) is a matter
configuration in a rotated duality frame rather than a further sector, and the light-like kernels form the only orbit fixed by the whole
duality circle.  This is the precise sense in which, on the kernel, matter and tachyon are one object in two frames and light is the
frame-invariant fix-sector.  Whether the duality angle is physical or a gauge freedom is not decided at this level.

\subsection{The compensator}

\begin{definition}[Involution and compensator]\label{def:compensator}
Let \(J_Bz:=i\bar z\), with complex conjugation of the components.  \(J_B\) is antilinear and \(J_B^2=1\).  The \emph{compensator} and its partner
are \(C:=\tfrac12(1+J_B)\) and \(A:=\tfrac12(1-J_B)\).
\end{definition}

In terms of fields, \(J_B(E+iB)=B+iE\) exchanges the electric and the magnetic field.  The compensator is the output of a two-arm
interferometer whose arms carry \(f\) and \(J_Bf\); its two ports carry \(Cf\) and \(Af\).

\begin{proposition}[Two ports]\label{prop:ports}
\(C\) and \(A\) are real-linear projections with \(C+A=1\) and \(CA=AC=0\).  With respect to the real inner product
\(\operatorname{Re}\langle Cf,Af\rangle=0\), so \(|Cf|^2+|Af|^2=|f|^2\).  Moreover \(A(f)=-i\,C(if)\).
\end{proposition}

\begin{proof}
The identities follow from \(J_B^2=1\).  Since \(|J_Bf|=|f|\),
\(\operatorname{Re}\langle Cf,Af\rangle=\tfrac14\operatorname{Re}\bigl(\langle J_Bf,f\rangle-\langle f,J_Bf\rangle\bigr)=0\).  Finally
\(C(if)=\tfrac12(if-iJ_Bf)=iA(f)\).
\end{proof}

The two ports are therefore real and imaginary parts of one field, not independent sectors.  Two examples show what the compensator does.
For the Coulomb field of a charge \(q\) (matter type), \(J_Bf\) is the field of a magnetic charge \(q\) (tachyon type), and \(Cf\), \(Af\) are the
fields of dyons with electric charge \(q/2\) and magnetic charge \(\pm q/2\).  For a plane light wave, \(J_B\) reverses the direction of
propagation and rotates the polarization by \(90^\circ\), as a phase-conjugating mirror followed by a polarization rotator would; \(Cf\) has
\(E\parallel B\) and is not a light wave.

\subsection{The doublet and the state disk}

Since \(J_B\) is antilinear, \(C\) has no complex matrix.  It becomes an ordinary matrix on the real doublet \(F=[\,E\mid B\,]\in\mathbb R^{3\times2}\).

\begin{proposition}[Doublet form]\label{prop:doublet}
Acting on \(F\) from the right, \(J_B\) is \(\sigma_x\), complex conjugation is \(\sigma_z\), and the duality rotation by \(\alpha\) is \(R_\alpha^{\mathsf T}\);
they generate \(O(2)\), and spatial rotations act from the left and commute with all of them.  The compensator is the orthogonal projection
\(C_{\rm HC}=\tfrac12(1+\sigma_x)\).  The Gram matrix \(G=F^{\mathsf T}F\) decomposes as
\begin{equation}\label{eq:gram}
G=u\,\mathbf 1+\tfrac12\bigl(\operatorname{Re}s\;\sigma_z+\operatorname{Im}s\;\sigma_x\bigr),\qquad u=\tfrac12(E^2+B^2),\qquad
\det G=|E\times B|^2 .
\end{equation}
\end{proposition}

\begin{proof}
\(i(E-iB)=B+iE\) exchanges the columns, \(E-iB\) changes the sign of the second, and \(e^{i\alpha}(E+iB)\) rotates them.  The entries of \(G\) are
\(E^2\), \(B^2\) and \(E\cdot B\).
\end{proof}

Put \(\rho=G/\operatorname{tr}G=\tfrac12(\mathbf 1+b_1\sigma_z+b_2\sigma_x)\) with \(b=(\operatorname{Re}s,\operatorname{Im}s)/(2u)\).  This is the Stokes formalism
for the doublet \((E,B)\) in place of two polarization components, and it gives a simple picture of the three sectors:

\begin{proposition}[State disk]\label{prop:disk}
\(|b|\le1\), with \(|b|=1\) exactly when \(E\parallel B\) and \(b=0\) exactly for a null field.  Matter-type fields have \(b_1>0\), tachyon-type
fields \(b_1<0\), light-like fields lie at the centre.  A duality rotation by \(\alpha\) rotates \(b\) by \(2\alpha\); \(J_B\) reflects \(b_1\); the quarter
turn \eqref{eq:quarterturn} reverses \(b\).  Under a boost \(s\) is invariant and \(u\) changes, so \(b\) moves along its own direction; it reaches
\(|b|=1\) in the frame with \(E\parallel B\) and tends to \(0\) as the boost grows.
\end{proposition}

\begin{proof}
\(u^2-\det G=\tfrac14|s|^2\) gives \(|b|^2=1-\det G/u^2\le1\).  A duality rotation conjugates \(G\) by \(R_\alpha\), which turns its traceless part by
\(2\alpha\).  The boost statement follows from the invariance of \(s\) and the growth of \(u\).
\end{proof}

Light is the unpolarized centre of the disk, matter and tachyon are its two fully polarized ends along one axis, and boosting any field
moves it towards light.  Both \(J_B\) and the quarter turn take matter-type to tachyon-type states; they differ in the second invariant
\(\operatorname{Im}s=2E\cdot B\), which \(J_B\) keeps and the quarter turn reverses.
